\chapter{Logical Design}

\section*{Volumes Table}

We consider a one-year observation window for volume estimation.

\begin{table}[H]
    \centering
    \resizebox{\textwidth}{!}{
    \begin{tabular}{|l|c|c|p{9.5cm}|}
    \hline
    \textbf{Concept} & \textbf{Type} & \textbf{Volume} & \textbf{Why} \\
    \hline
    Depot & E & 10 & By specification. \\
    \hline
    Setup Team & E & 100 & By specification. \\
    \hline
    Member & E & 1000 & By specification: each setup team can have at most 10 members, so $100 \cdot 10 = 1000$. \\
    \hline
    Booking & E & 50000 & 1:1 with Place. \\
    \hline
    Special Booking & E & 1000 & Assumption: special rentals are much more occasional than the other booking types. \\
    \hline
    One-Time Booking & E & 40000 & Assumption: one-time bookings are the most frequent booking type. \\
    \hline
    Recurrent Booking & E & 9000 & Assumption: recurrent bookings include seasonal contracts and represent the remaining recurring workload. \\
    \hline
    Event Location & E & 1000 & By specification. \\
    \hline
    Customer & E & 500 & Assumption: each customer has on average at least 2 event locations, so $1000 / 2 = 500$. \\
    \hline
    Individual Customer & E & 400 & Assumption: 4/5 of customers are individuals. \\
    \hline
    Company Customer & E & 100 & Assumption: 1/5 of customers are companies. \\
    \hline
    Employ & R & 100 & 1:1 with Setup Team. \\
    \hline
    Works In & R & 1000 & 1:1 with Member. \\
    \hline
    Handle & R & 50000 & 1:1 with Booking. \\
    \hline
    Delivered To & R & 50000 & 1:1 with Booking. \\
    \hline
    Has & R & 1000 & 1:1 with Event Location. \\
    \hline
    Place & R & 50000 & Assumption: each customer places 100 bookings, so $500 \cdot 100 = 50{,}000$. \\
    \hline
    \end{tabular}
    }
\end{table}

\section*{Operations Table}

The operation set is taken directly from the track.

\begin{table}[H]
    \centering
    \begin{tabular}{|c|c|p{10cm}|c|}
    \hline
    \textbf{Op.} & \textbf{Type} & \textbf{Description} & \textbf{Frequency} \\
    \hline
    1 & I & Enter all data related to a new customer & 10/day \\
    \hline
    2 & I & Record a new booking & 300/day \\
    \hline
    3 & I & Register a new event location & 50/day \\
    \hline
    4 & I & View the team that handled setups at a specific event location & 20/day \\
    \hline
    5 & I & Print event locations sorted by number of bookings (descending) & 5/day \\
    \hline
    \end{tabular}
\end{table}

\section*{Operations Analysis}

\subsection*{Operation 1: Enter all data related to a new customer}

\begin{table}[H]
    \centering
    \begin{tabular}{|l|c|c|c|p{6.3cm}|}
    \hline
    \textbf{Concept} & \textbf{Type} & \textbf{ACC} & \textbf{ACC Type} & \textbf{Why} \\
    \hline
    Customer & E & 1 & W & Create the customer. \\
    \hline
    Event Location & E & 2 & W & Create the two event locations. \\
    \hline
    Customer & E & 1 & R & Read the ID of the newly created customer. \\
    \hline
    Event Location & E & 2 & R & Read the two event locations for the customer. \\
    \hline
    Has & R & 2 & W & Create the two ownership relationships. \\
    \hline
    \end{tabular}
\end{table}

\textbf{Operation cost}: $(2 \cdot (1 + 2 + 2) + (1 + 2)) \cdot 10 = 130$ accesses/day.

\subsection*{Operation 2: Record a new booking}

For creating a new booking, we usually start from the customer ID and the code of the selected event location. Then, based on that location, we identify the setup team that covers the corresponding municipality/area.

Access \textit{without} redundancy:

\begin{table}[H]
    \centering
    \begin{tabular}{|l|c|c|c|p{6.3cm}|}
    \hline
    \textbf{Concept} & \textbf{Type} & \textbf{ACC} & \textbf{ACC Type} & \textbf{Why} \\
    \hline
    Booking & E & 1 & W & Create the booking. \\
    \hline
    Delivered To & R & 1 & W & Create the relationship with the selected event location. \\
    \hline
    Place & R & 1 & W & Create the relationship with the customer. \\
    \hline
    Event Location & E & 1 & R & Read the city/province of the selected event location to identify the serving depot. \\
    \hline
    Depot & E & 10 & R & Find the depot that covers the area of the selected event location. \\
    \hline
    Employ & R & 10 & R & Find the candidate setup teams linked to the selected depot. \\
    \hline
    Handle & R & 1 & W & Create the relationship between booking and selected setup team. \\
    \hline
    \end{tabular}
\end{table}

\textbf{Operation cost} (\textit{without redundancy}): $29 \cdot 300 = 8700$ accesses/day.

Access \textit{with} redundancy:

\begin{table}[H]
    \centering
    \begin{tabular}{|l|c|c|c|p{6.3cm}|}
    \hline
    \textbf{Concept} & \textbf{Type} & \textbf{ACC} & \textbf{ACC Type} & \textbf{Why} \\
    \hline
    Booking & E & 1 & W & Create the booking. \\
    \hline
    Delivered To & R & 1 & W & Create the relationship with the selected event location. \\
    \hline
    Place & R & 1 & W & Create the relationship with the customer. \\
    \hline
    Event Location & E & 1 & R & Read the city/province of the selected event location to identify the serving depot. \\
    \hline
    Depot & E & 10 & R & Find the depot that covers the area of the selected event location. \\
    \hline
    Employ & R & 10 & R & Find the candidate setup teams linked to the selected depot. \\
    \hline
    Setup Team & E & 1 & R & Read the \texttt{n. Installations} of the selected setup team. \\
    \hline
    Setup Team & E & 1 & W & Update the \texttt{n. Installations} for the selected setup team. \\
    \hline
    Handle & R & 1 & W & Create the relationship between booking and selected setup team. \\
    \hline
    \end{tabular}
\end{table}

\textbf{Operation cost} (\textit{with redundancy}): $32 \cdot 300 = 9600$ accesses/day.

\subsection*{Operation 3: Register a new event location}

For registering a new event location, we usually already have the customer ID, since the customer must be authenticated before adding a location.

\begin{table}[H]
    \centering
    \begin{tabular}{|l|c|c|c|p{6.3cm}|}
    \hline
    \textbf{Concept} & \textbf{Type} & \textbf{ACC} & \textbf{ACC Type} & \textbf{Why} \\
    \hline
    Event Location & E & 1 & W & Create the event location. \\
    \hline
    Has & R & 1 & W & Create the relationship with the customer. \\
    \hline
    \end{tabular}
\end{table}

\textbf{Operation cost}: $4 \cdot 50 = 200$ accesses/day.

\subsection*{Operation 4: View the team that handled setups at a specific event location}

\begin{table}[H]
    \centering
    \begin{tabular}{|l|c|c|c|p{6.3cm}|}
    \hline
    \textbf{Concept} & \textbf{Type} & \textbf{ACC} & \textbf{ACC Type} & \textbf{Why} \\
    \hline
    Event Location & E & 1 & R & Read the ID of the selected event location. \\
    \hline
    Delivered To & R & 50 & R & Average of 50 bookings per location: $50000 / 1000 = 50$. \\
    \hline
    Booking & E & 50 & R & 1:1 with Delivered To, one booking for each delivered-to instance. \\
    \hline
    Handle & R & 50 & R & 1:1 with Booking, retrieve the team relation for each booking. \\
    \hline
    \end{tabular}
\end{table}

\textbf{Operation cost}: $151 \cdot 20 = 3020$ accesses/day.

\subsection*{Operation 5: Print event locations sorted by number of bookings}

\begin{table}[H]
    \centering
    \begin{tabular}{|l|c|c|c|p{6.3cm}|}
    \hline
    \textbf{Concept} & \textbf{Type} & \textbf{ACC} & \textbf{ACC Type} & \textbf{Why} \\
    \hline
    Delivered To & R & 50000 & R & Read all relationship instances to count occurrences of each event location ID. \\
    \hline
    Event Location & E & 1000 & R & Access and print the actual event location data. \\
    \hline
    \end{tabular}
\end{table}

\textbf{Operation cost}: $(50000 + 1000) \cdot 5 = 255000$ accesses/day.

\section*{Design Considerations}

In this section, redundancy and generalization situations are analyzed.

\subsection*{Setup Team n. Installations}

The attribute \texttt{n. Installations} in \texttt{Setup Team} is removed from the logical baseline as a stored redundancy. The comparison on Operation 2 shows that maintaining this counter increases write cost (additional read/write on setup team), while none of the required operations explicitly needs that value as direct output.

\subsection*{Relation Place (Customer-Booking)}

The relation \texttt{Place} is redundant, because customer information is still reachable through \texttt{Has} (Customer-Event Location) together with \texttt{Delivered To} (Booking-Event Location).

\textbf{Operation 2 with \texttt{Place}} (original table):

\begin{table}[H]
    \centering
    \begin{tabular}{|l|c|c|c|p{6.3cm}|}
    \hline
    \textbf{Concept} & \textbf{Type} & \textbf{ACC} & \textbf{ACC Type} & \textbf{Why} \\
    \hline
    Booking & E & 1 & W & Create the booking. \\
    \hline
    Delivered To & R & 1 & W & Create the relationship with the selected event location. \\
    \hline
    Place & R & 1 & W & Create the relationship with the customer. \\
    \hline
    Event Location & E & 1 & R & Read the city/province of the selected event location to identify the serving depot. \\
    \hline
    Depot & E & 10 & R & Find the depot that covers the area of the selected event location. \\
    \hline
    Employ & R & 10 & R & Find the candidate setup teams linked to the selected depot. \\
    \hline
    Handle & R & 1 & W & Create the relationship between booking and selected setup team. \\
    \hline
    \end{tabular}
\end{table}

\textbf{Operation 2 with \texttt{Has} replacing \texttt{Place}}:

\begin{table}[H]
    \centering
    \begin{tabular}{|l|c|c|c|p{6.3cm}|}
    \hline
    \textbf{Concept} & \textbf{Type} & \textbf{ACC} & \textbf{ACC Type} & \textbf{Why} \\
    \hline
    Booking & E & 1 & W & Create the booking. \\
    \hline
    Delivered To & R & 1 & W & Create the relationship with the selected event location. \\
    \hline
    \textbf{Has} & \textbf{R} & \textbf{1} & \textbf{R} & \textbf{Check customer ownership of the selected event location (replacement row).} \\
    \hline
    Event Location & E & 1 & R & Read the city/province of the selected event location to identify the serving depot. \\
    \hline
    Depot & E & 10 & R & Find the depot that covers the area of the selected event location. \\
    \hline
    Employ & R & 10 & R & Find the candidate setup teams linked to the selected depot. \\
    \hline
    Handle & R & 1 & W & Create the relationship between booking and selected setup team. \\
    \hline
    \end{tabular}
\end{table}

\[
\text{With Place: } (2 \cdot 4 + 21) \cdot 300 = 29 \cdot 300 = 8700\ \text{accesses/day}
\]
\[
\text{With Has replacement: } (2 \cdot 3 + 22) \cdot 300 = 28 \cdot 300 = 8400\ \text{accesses/day}
\]

\noindent The replacement preserves functionality and slightly reduces weighted accesses, so \texttt{Place} is removed from the logical restructuring.

\subsection*{Partitioning and Merging Choices}

For \texttt{Booking}, the specialization (Special / One-Time / Recurrent) is collapsed into a single entity. The entity keeps standard booking attributes and introduces:

\begin{itemize}
    \item a boolean \texttt{IsSpecial};
    \item an \texttt{Interval} attribute for recurrent behavior.
\end{itemize}

Interpretation rule:

\begin{itemize}
    \item if \texttt{IsSpecial = TRUE}, the booking is a special rental;
    \item if \texttt{Interval} is greater than zero, the booking is recurrent.
\end{itemize}

For \texttt{Customer}, specialization is maintained, because \texttt{Individual} and \texttt{Company} have distinct, non-overlapping attributes.

\subsection*{Event Location Identifier After Place Removal}

After removing \texttt{Place}, \texttt{Event Location} is treated with an autonomous unique code (no longer customer-scoped). This keeps the \texttt{Booking}-\texttt{Event Location} connection independent and supports the replacement of \texttt{Place} with the \texttt{Delivered To} + \texttt{Has} navigation path.

\section*{UML Schema}

The following UML schema summarizes the final logical structure adopted in this chapter.

\begin{figure}[H]
    \centering
    \safeincludegraphics[width=0.92\textwidth]{images/UMLFINAL.png}
    \caption{UML Schema}
\end{figure}







